\section{Aplicación práctica del modelado de secuencias: modelado del lenguaje}

\subsection{La tarea de Next Word Prediction}

El modelado del lenguaje (Language Modeling) es una de las aplicaciones más representativas e importantes del modelado de secuencias. Su tarea central es extremadamente simple: \textbf{dada una secuencia de palabras, predecir la siguiente palabra más probable}.

\begin{figure}[H]
\centering
\includegraphics[width=0.85\textwidth]{slides-images/slide-35.jpg}
\caption{Predecir la siguiente palabra: la tarea central del modelado del lenguaje\protect\footnotemark}
\end{figure}
\footnotetext{Diapositiva 35.}

Por ejemplo, dada la oración ``This morning I took my cat for a'', el modelo necesita predecir que la siguiente palabra es ``walk''.

\begin{importantbox}{Next Word Prediction es la base de los LLM}
Los grandes modelos de lenguaje actuales (como la serie GPT) están, en esencia, ejecutando esta tarea de Next Word Prediction. Al entrenarse sobre enormes cantidades de texto para predecir la siguiente palabra, el modelo aprende la estructura y la semántica del lenguaje.
\end{importantbox}

\subsection{Cómo representar el lenguaje mediante números}

Las redes neuronales solo pueden procesar datos numéricos, por lo que necesitamos convertir las palabras en vectores numéricos. Este proceso consta de dos pasos:

\textbf{Paso 1: Construir el vocabulario (Vocabulary)}

Definir un vocabulario finito que contenga todas las palabras posibles y asignar a cada palabra un número de índice único:
\begin{itemize}
    \item ``a'' $\rightarrow$ 1
    \item ``cat'' $\rightarrow$ 2
    \item ``the'' $\rightarrow$ 3
    \item $\ldots$
\end{itemize}

\textbf{Paso 2: Embedding}

Mapear los índices a vectores de dimensión fija. El método más simple es \textbf{One-Hot Encoding}:

$$\text{``cat''} \rightarrow [0, 1, 0, 0, \ldots, 0]$$

\begin{warningbox}{Las limitaciones del One-Hot Encoding}
El One-Hot Encoding tiene dos defectos graves:
\begin{enumerate}
    \item \textbf{La maldición de la dimensionalidad}: la dimensión del vector es igual al tamaño del vocabulario; para un vocabulario con decenas o incluso cientos de miles de palabras, esta representación resulta extremadamente dispersa e ineficiente
    \item \textbf{Ausencia de información semántica}: la ``distancia'' entre todas las palabras es igual; la distancia entre ``cat'' y ``dog'' es la misma que entre ``cat'' y ``airplane'', por lo que no puede capturar la similitud semántica
\end{enumerate}
\end{warningbox}

Un método mejor es usar \textbf{Learned Embedding} (embedding aprendido): mediante entrenamiento, la red aprende de manera automática una representación vectorial de baja dimensión y densa, de modo que las palabras semánticamente cercanas queden a menor distancia en el espacio vectorial.

\begin{figure}[H]
\centering
\includegraphics[width=0.75\textwidth]{slides-images/slide-38.jpg}
\caption{Learned Embedding: las palabras semánticamente cercanas quedan a menor distancia en el espacio vectorial\protect\footnotemark}
\end{figure}
\footnotetext{Diapositiva 38.}

\begin{knowledgebox}{La naturaleza del embedding}
Un embedding es, en esencia, una tabla de consulta (Lookup Table) que puede aprenderse. Mapea símbolos discretos (como palabras o caracteres) a un espacio vectorial continuo. Después del entrenamiento, las palabras semánticamente similares (como ``king'' y ``queen'') quedan cercanas entre sí en el espacio de embedding, mientras que las palabras con significados distintos quedan más alejadas. Esta propiedad resulta fundamental para las tareas posteriores de modelado de secuencias.
\end{knowledgebox}

\subsection{Dependencias de corto y largo plazo en las secuencias}

En el lenguaje natural existen relaciones de dependencia de distinto alcance:

\begin{itemize}
    \item \textbf{Dependencia de corto plazo}: ``The clouds are in the \underline{\hspace{1cm}}'' $\rightarrow$ ``sky''. La respuesta depende solo de las palabras inmediatamente anteriores.
    \item \textbf{Dependencia de largo plazo}: ``I grew up in France, $\ldots$ and I speak fluent \underline{\hspace{1cm}}'' $\rightarrow$ ``French''. La respuesta depende de un contexto muy anterior, ``France''.
\end{itemize}

Las dependencias de largo plazo plantean un desafío severo para las RNN, lo cual da paso directamente al tema de la siguiente sección.

\subsection{Resumen de la sección}

La tarea central del modelado del lenguaje —predecir la siguiente palabra— parece simple pero es extremadamente poderosa: es la base de entrenamiento de todos los grandes modelos de lenguaje modernos. Representar el lenguaje mediante vectores numéricos (embedding) es el primer paso para lograr esta tarea, mientras que manejar las dependencias de corto y largo plazo en la secuencia constituye el desafío central.

\newpage
